\documentclass[aspectratio=169]{beamer}

\usetheme{Madrid}
\usepackage{booktabs}
\usepackage{graphicx}
\setbeamertemplate{navigation symbols}{}
\setbeamertemplate{caption}{\raggedright\insertcaption\par}

\title[Refining the LULC map]{From 4 fixed classes to a user-grown hierarchy}
\subtitle{Splitting and adding LULC classes, live, at 10\,m}
\author[Salil Gujar]{Salil Gujar \\ \footnotesize M.Tech CSE \quad Entry No.\ 2025MCS2106}
\date{June 2026}

\begin{document}

\frame{\titlepage}

% ------------------------------------------------------------------
\begin{frame}{This week}
\begin{itemize}
  \item A 4-class base map (greenery, water, built-up, barren) from Alpha Earth, at 10\,m.
  \item This week: let the user \textbf{grow the classes into a hierarchy} on top of that
  base --- two operations, at \emph{any} level:
  \begin{itemize}
    \item \textbf{SPLIT} a class into children (greenery $\to$ crops / trees / shrubs);
    \item \textbf{ADD} a new class under a node (barren $\to$ +\,mining).
  \end{itemize}
  \item Mark or upload a few example polygons, \textbf{retrain on the fly}, and the 10\,m
  map updates with the new sub-classes.
\end{itemize}
\end{frame}

% ------------------------------------------------------------------
\begin{frame}{The class tree}
\begin{itemize}
  \item One JSON tree; each node carries \texttt{class, name, parent, color, children,
  source} (where its training pixels come from).
  \item \textbf{Every non-leaf node owns one small classifier} over just its children.
  Inference walks \textbf{root $\to$ leaf}: base map says \emph{greenery}, the greenery
  classifier says \emph{crops}, and so on --- recursively, to any depth.
  \item SPLIT and ADD are the \emph{same} machinery: ADD onto a leaf is a SPLIT into an
  auto ``residual'' (keeps the parent's identity) plus the new class.
  \item Per-node, not one giant model: a new class refits \textbf{one} classifier on
  \textbf{its} data in seconds, leaving the base map untouched.
\end{itemize}
\end{frame}

% ------------------------------------------------------------------
\begin{frame}{Demo 1 --- SPLIT greenery $\to$ crops / trees / shrubs}
\begin{columns}[T]
\begin{column}{0.52\textwidth}
\footnotesize
\textbf{Data:} crops \& trees from FarmForest \emph{expert} polygons (sampled through the
tool); shrubs from ESA WorldCover (we couldn't get expert shrub polygons this week).
\vspace{0.4em}
\centering
\begin{tabular}{l c c c}
\toprule
class & P & R & F1 \\
\midrule
crops  & 0.987 & 0.943 & 0.965 \\
trees  & 0.997 & 0.986 & 0.991 \\
shrubs & 0.548 & 0.947 & 0.695 \\
\midrule
\multicolumn{4}{c}{accuracy \textbf{0.963} (3{,}014 held-out px)} \\
\bottomrule
\end{tabular}
\flushleft\vspace{0.4em}
Polygon-held-out split (no pixel leakage). Shrubs is the weak class --- few, noisier
WorldCover labels.
\end{column}
\begin{column}{0.48\textwidth}
\centering
\includegraphics[width=0.49\linewidth]{figures/greenery_before.png}\hfill
\includegraphics[width=0.49\linewidth]{figures/greenery_after.png}\\
\footnotesize before: greenery \quad$\to$\quad after: crops / trees
\end{column}
\end{columns}
\end{frame}

% ------------------------------------------------------------------
\begin{frame}{Demo 2 --- ADD mining under barren}
\begin{columns}[T]
\begin{column}{0.52\textwidth}
\footnotesize
\textbf{Data:} mining positives from the provided mining polygons
(\texttt{mining\_polygons\_india.gpkg}); the residual ``Barren'' from existing barren
pixels --- so ADD reduces to a 2-way split.
\vspace{0.4em}
\centering
\begin{tabular}{l c c c}
\toprule
class & P & R & F1 \\
\midrule
mining        & 0.860 & 0.886 & 0.873 \\
barren (rest) & 0.850 & 0.817 & 0.833 \\
\midrule
\multicolumn{4}{c}{accuracy \textbf{0.856} (4{,}294 held-out px)} \\
\bottomrule
\end{tabular}
\flushleft\vspace{0.4em}
The map gains a \emph{mining} class with no change to the inference code --- it
composites like any other split.
\end{column}
\begin{column}{0.48\textwidth}
\centering
\includegraphics[width=0.49\linewidth]{figures/mining_before.png}\hfill
\includegraphics[width=0.49\linewidth]{figures/mining_after.png}\\
\footnotesize before: barren \quad$\to$\quad after: mining / barren
\end{column}
\end{columns}
\end{frame}

% ------------------------------------------------------------------
\begin{frame}{Any level of the hierarchy}
\begin{itemize}
  \item \textbf{Splits nest:} a leaf can itself be split one level deeper, and inference
  recurses into it unchanged --- the tree can grow to any depth.
  \item \textbf{Root level} too: adding a brand-new \emph{base} class retrains the pooled
  base model (Alpha Earth + WorldCover), preserving its calibration.
  \item Validated end-to-end: a 3-deep tree (root $\to$ greenery $\to$ trees $\to$ two
  sub-classes) renders correctly --- the server-side map matched the Python model
  \textbf{36/36} at sampled pixels.
\end{itemize}
\end{frame}

% ------------------------------------------------------------------
\begin{frame}{One classifier per node, or one big multiclass?}
We compared both on the greenery leaves (shared, polygon-held-out test, 15{,}573 px):
\vspace{0.4em}
\centering
\small
\begin{tabular}{l c c c c c}
\toprule
\textbf{topology} & \textbf{acc} & \textbf{macro-F1} & crops F1 & trees F1 & shrubs F1 \\
\midrule
Flat (one 6-way model)        & \textbf{0.885} & \textbf{0.777} & 0.842 & 0.938 & 0.158 \\
Hierarchical (base $\to$ split) & 0.862 & 0.750 & 0.779 & 0.912 & 0.133 \\
\bottomrule
\end{tabular}
\vspace{0.8em}
\flushleft
\begin{itemize}\small
  \item Flat is $\approx$2 pts better offline (no cascading base errors), but
  \textbf{per-node wins for the live tool}: each edit refits one tiny model on its own
  data, the base map stays stable, and errors stay local.
  \item Shrubs is weak in \emph{both} --- a data gap (WorldCover-only), not a topology
  one.
\end{itemize}
\end{frame}

% ------------------------------------------------------------------
\begin{frame}{10\,m, server-side}
\begin{itemize}
  \item The base map and every split are \textbf{linear} (\texttt{StandardScaler $\to$
  LinearSVC}), so we reproduce them \emph{exactly} as band math on the Alpha Earth image
  \textbf{inside Earth Engine} and classify at native 10\,m.
  \item One classified PNG per area (no per-point sampling, no payload blow-up); the
  greenery split composites on top, then mining, then deeper splits.
  \item Correctness check: the Earth-Engine label image matched the trained sklearn
  pipeline \textbf{36/36} at sampled points (base + both splits).
\end{itemize}
\end{frame}

% ------------------------------------------------------------------
\begin{frame}{The interactive loop}
\begin{itemize}
  \item Browser tool (FastAPI + Leaflet): a \textbf{hierarchy tree} you click to select a
  class.
  \item \textbf{Add examples} by drawing on the map or uploading GeoJSON / KML; positive
  = ``this is X'' (relabel), negative = ``not X'' (hard-negative).
  \item \textbf{Split / Add / Retrain} a node from the sidebar --- training runs, metrics
  print, and the 10\,m raster re-renders with the new classes.
\end{itemize}
\vspace{0.6em}
\textbf{Honest limits:} numbers are agreement vs expert / WorldCover labels, not field
truth; shrubs leans on WorldCover (its weak class); root-level retrain is heavier and
touches the shared base map.
\end{frame}

% ------------------------------------------------------------------
\begin{frame}{How accurate is the base map?}
Two honest tests of the 4-class base (Realistic --- Alpha Earth $+$ WorldCover, 10\,m):
\vspace{0.4em}
\centering
\small
\begin{tabular}{l c c}
\toprule
\textbf{test} & \textbf{accuracy} & \textbf{macro-F1} \\
\midrule
Balanced expert hold-out (classes even) & \textbf{0.89} & \textbf{0.88} \\
Random real-India ($\approx$89\% greenery) & 0.80 & --- \\
\bottomrule
\end{tabular}
\vspace{0.8em}
\flushleft
\begin{itemize}\small
  \item Strong per-class skill: \textbf{0.89} on the balanced test. On random real ground,
  greenery is near-perfect (precision 0.99); the 0.80 is held back only by the \emph{rare}
  classes (barren / water), which are genuinely hard there.
  \item vs the alternative: Realistic beats Detailed (AE $+$ Tessera) on random India
  (0.80 vs 0.76) --- Tessera doesn't earn its place under the real prior.
  \item Both modes kept: Realistic is the default 10\,m map; Detailed for rare-class detail.
\end{itemize}
\end{frame}

% ------------------------------------------------------------------
\begin{frame}
\centering
\Large Thank you 
\end{frame}

\end{document}
